\chapter{2026 春季 QE-AI 真题解答}
\label{ch:2026spring}

\section{A. Machine Learning Theory}

\begin{examproblem}[A.MQ1-MQ6：选择题]
agnostic PAC、Sauer 引理、ReLU 学习理论、ERM、No Free Lunch 和 kernel 的选择题。

\source{官方卷：2026 Spring, p.~1--2（已按 PDF 页面视觉复核）}
\officialenglish{\textbf{MQ1} [0.5 pts] In the agnostic PAC learning setting, what quantity does a learning algorithm aim to compete with? (a) The Bayes optimal classifier. (b) The hypothesis with zero training error. (c) The best hypothesis in the class \(\mathcal H\). (d) The hypothesis minimising validation error.

\textbf{MQ2} [0.5 pts] Let \(\mathcal H\) be a binary hypothesis class with VC dimension \(d<\infty\). Which statement about its growth function \(\tau_{\mathcal H}(m)\) is necessarily true? (a) \(\tau_{\mathcal H}(m)=2^m\) for all \(m\). (b) \(\tau_{\mathcal H}(m)=O(m^d)\) for all \(m\). (c) \(\tau_{\mathcal H}(m)\le\sum_{i=0}^{d}\binom mi\) for all \(m\). (d) \(\tau_{\mathcal H}(m)\) is constant for \(m>d\).

\textbf{MQ3} Which statement about ReLU neural networks is correct from a learning-theoretic perspective? (a) ReLU networks represent smooth functions on \(\mathbb R^d\). (b) Increasing depth always decreases the VC dimension. (c) ReLU networks compute piecewise linear functions whose complexity depends on depth and width. (d) Universal approximation of ReLU networks implies PAC learnability.

\textbf{MQ4} Which statement about empirical risk minimisation (ERM) is correct? (a) ERM is guaranteed to be consistent for any hypothesis class. (b) ERM always finds a hypothesis with minimum true risk. (c) If \(\mathcal H\) has finite VC dimension, ERM is PAC learnable in the realisable setting. (d) ERM is PAC learnable only if \(\mathcal H\) is finite.

\textbf{MQ5} Which statement is a correct consequence of the No Free Lunch theorem? (a) Restricting the hypothesis class always improves generalization. (b) There exists a universally optimal learning algorithm. (c) Inductive bias is necessary to obtain non-trivial learning guarantees. (d) Random guessing is optimal for all supervised learning problems.

\textbf{MQ6} Let \(K:X\times X\to\mathbb R\) be symmetric. Which statement is correct? (a) \(K\) is a valid kernel iff \(K(x,x)\ge0\) for all \(x\). (b) iff it has a finite-dimensional feature map. (c) If \(K\) is positive semidefinite, there exist a possibly infinite-dimensional Hilbert space \(\mathcal H\) and \(\psi:X\to\mathcal H\) such that \(K(x,x')=\langle\psi(x),\psi(x')\rangle_{\mathcal H}\). (d) Every kernel has a unique feature map.}
\chinesetranslation{六道单选题依次考查：agnostic PAC 的比较对象、有限 VC 维类的增长函数、ReLU 的分段线性表达、ERM 的可学习性、No Free Lunch 的归纳偏置结论，以及正半定核的 Moore--Aronszajn 表示。每题四选一。MQ2 的 (c) 中二项式系数是 PDF 原图可见而文本抽取损坏的部分。}
\end{examproblem}

\solution
答案为

\[
  \mathrm{MQ1}(c),\quad
  \mathrm{MQ2}(c),\quad
  \mathrm{MQ3}(c),\quad
  \mathrm{MQ4}(c),\quad
  \mathrm{MQ5}(c),\quad
  \mathrm{MQ6}(c).
\]

其中 MQ2 的官方 PDF 抽取文本把二项式符号显示得不完整；按 \conceptref{concept:vc-sauer}{Sauer 引理}，正确选项应为

\[
  \tau_{\Hcal}(m)\le \sum_{i=0}^{d}\binom{m}{i}.
\]

MQ6 是 Moore-Aronszajn 定理：symmetric positive semidefinite \conceptref{concept:svm-kernel}{kernel} 对应某个可能无限维 Hilbert space 的 feature map，但 feature map 不唯一。

\begin{examproblem}[A1：经验损失均方偏差]
设 \(p=L_{\Dcal,f}(h)\)，证明

\[
  \E_{S\sim \Dcal^m}\left[(L_S(h)-p)^2\right]
  =
  \frac{p(1-p)}{m}
  \le \frac{1}{4m}.
\]

\source{官方卷：2026 Spring, p.~2}
\officialenglish{Let \(\mathcal H\) be a hypothesis class of binary classifiers \(h:X\to\{0,1\}\). Let \(D\) be an unknown distribution over \(X\), and let \(f\in\mathcal H\) be the target hypothesis. For fixed \(h\in\mathcal H\), define \(L_S(h)=m^{-1}\sum_{i=1}^m\mathbf 1[h(x_i)\ne f(x_i)]\) on \(S=\{x_1,\ldots,x_m\}\sim D^m\). (a) Let \(p:=L_{D,f}(h)\). Show \(\mathbb E_{S\sim D^m}[(L_S(h)-p)^2]=p(1-p)/m\). (b) Deduce that this is at most \(1/(4m)\), hence decreases at rate \(O(1/m)\).}
\chinesetranslation{设二分类假设类 \(\mathcal H\)、未知分布 \(D\)、目标假设 \(f\in\mathcal H\)，固定 \(h\)。对 \(m\) 个独立样本的经验损失，令 \(p=L_{D,f}(h)\)：(a) 证明其相对 \(p\) 的均方偏差为 \(p(1-p)/m\)；(b) 推出其不超过 \(1/(4m)\)，因而按 \(O(1/m)\) 衰减。}

\end{examproblem}

\solution
与 2025 春 A1 相同。令 \(Z_i=\mathbf{1}[h(x_i)\ne f(x_i)]\)，则 \(Z_i\sim\Bern(p)\) 且 \(L_S=m^{-1}\sum_iZ_i\)。因为 \(\E L_S=p\)，

\[
  \E[(L_S-p)^2]=\Var(L_S)=\frac{p(1-p)}{m}.
\]

又 \(p(1-p)\le 1/4\)，所以

\[
  \E[(L_S-p)^2]\le \frac{1}{4m}.
\]

\begin{examproblem}[A2：强凸函数的唯一全局极小点]
若 \(f:\R^d\to\R\) 二阶连续可微且 \(\lambda\)-strongly convex，证明其 global minimizer 唯一。

\source{官方卷：2026 Spring, p.~2}
\officialenglish{Let \(f:\mathbb R^d\to\mathbb R\) be twice continuously differentiable and \(\lambda\)-strongly convex. Prove that \(f\) admits a unique global minimiser.}
\chinesetranslation{设 \(f:\mathbb R^d\to\mathbb R\) 二阶连续可微且 \(\lambda\)-强凸。证明 \(f\) 存在唯一的全局极小点。}
\end{examproblem}

\solution
强凸性意味着对 \(x\ne y\) 和 \(t\in(0,1)\)，

\[
  f(tx+(1-t)y)
  <
  tf(x)+(1-t)f(y).
\]

若存在两个不同 global minimizers \(x^*,y^*\)，且最小值为 \(m\)，则

\[
  f(tx^*+(1-t)y^*)
  <
  tm+(1-t)m=m,
\]

与 \(m\) 是全局最小值矛盾。因此 minimizer 至多一个。若题设隐含 \(f\) coercive 或 minimizer 存在，则该 minimizer 唯一。对 \(\lambda\)-strongly convex 的二阶函数，Hessian 满足 \(\nabla^2f(x)\succeq \lambda I\)，通常也保证 \(f(x)\to\infty\) 当 \(\|x\|\to\infty\)，从而存在全局 minimizer。

\begin{examproblem}[A3：一维区间加右射线类]

\[
  h_{a,b,c}(x)=\mathbf{1}_{[a,b]}(x)\vee \mathbf{1}_{[c,\infty)}(x),
  \qquad a\le b<c.
\]

给出增长函数上界并求 VC 维。

\source{官方卷：2026 Spring, p.~3}
\officialenglish{Let \(X=\mathbb R\) and consider \(\mathcal H=\{h_{a,b,c}(x)=\mathbf 1_{[a,b]}(x)\vee\mathbf 1_{[c,\infty)}(x):a\le b<c,\ a,b,c\in\mathbb R\}\), where \(\vee\) denotes logical OR. (a) For \(n\) distinct points in \(\mathbb R\), give an upper bound on \(s(\mathcal H,n)\). (b) Determine \(\operatorname{VCdim}(\mathcal H)\).}
\chinesetranslation{在实线上考虑“一个闭区间的指示函数与一个右射线指示函数取逻辑或”的假设类，参数满足 \(a\le b<c\)。(a) 对 \(n\) 个互异点给出增长函数 \(s(\mathcal H,n)\) 的上界；(b) 求该类的 VC 维。}
\end{examproblem}

\solution
对 \(n\) 个有序点 \(x_1<\cdots<x_n\)，该类产生的正类集合形状是：

\begin{itemize}
  \item 空集；
  \item 一个连续正块 \([i,j]\)；
  \item 两个正块，其中第二个正块必须一直延伸到末尾，即 \([i,j]\cup[k,n]\)，且 \(i\le j<k-1\)。
\end{itemize}

因此不同 dichotomies 数为

\[
  s(\Hcal,n)
  \le
  1+\frac{n(n+1)}{2}
  +
  \sum_{j=1}^{n-2}j(n-j-1).
\]

计算求和：

\[
  \sum_{j=1}^{n-2}j(n-j-1)
  =
  \frac{n(n-1)(n-2)}{6}.
\]

所以

\[
  s(\Hcal,n)
  \le
  1+\frac{n(n+1)}{2}+\frac{n(n-1)(n-2)}{6}.
\]

当 \(n=3\) 时该数为 \(8=2^3\)，且三点确实可 shatter。例如模式 \(101\) 用第一个点的小区间加第三点开始的右射线；模式 \(010\) 用中间点的小区间且右射线放在第三点右侧。

当 \(n=4\) 时上界为

\[
  1+10+4=15<16,
\]

不能 shatter 四点。因此

\[
  \operatorname{VCdim}(\Hcal)=3.
\]

\begin{examproblem}[A4：多类 margin 构造]
样本 \((x,y)\) 在其类别中心 \(\mu_y\) 半径 \(r\) 球内，且不同中心距离至少 \(4r\)，\(r\ge1\)。用题给 multivector feature 和

\[
  w_i=(\mu_i,-\|\mu_i\|^2/2)
\]

证明存在 \(w\) 使所有样本 multiclass hinge loss 为零。
\source{官方卷：2026 Spring, p.~3}
\officialenglish{Consider examples \(S\subset\mathbb R^n\times[k]\). There are centres \(\mu_1,\ldots,\mu_k\in\mathbb R^n\) such that \(\|x-\mu_y\|\le r\) for each \((x,y)\in S\), where \(r\ge1\), and \(\|\mu_i-\mu_j\|\ge4r\) for \(i\ne j\). Concatenate every instance with a bias coordinate \(1\), and use the multivector construction \(\Psi(x,y)\in\mathbb R^{k(n+1)}\) with \((x,1)\) in the \(y\)-th block. Show that there is \(w\in\mathbb R^{k(n+1)}\) such that \(\langle w,\Psi(x,y)\rangle>0\) for every \((x,y)\in S\). Hint: take blocks \(w_i=[\mu_i,-\|\mu_i\|^2/2]\).}
\chinesetranslation{给定间隔至少 \(4r\) 的 \(k\) 个类别中心，每个样本到其真实类别中心的距离不超过 \(r\)。将 \((x,1)\) 放入对应类别的块向量 \(\Psi(x,y)\) 中，证明存在一个多向量分类器 \(w\)，使所有正确标注样本都有正间隔；提示取第 \(i\) 块为 \([\mu_i,-\|\mu_i\|^2/2]\)。}
\end{examproblem}

\solution
对类别 \(i\)，score 为

\[
  s_i(x)=\langle \mu_i,x\rangle-\frac12\|\mu_i\|^2.
\]

若样本真实类别为 \(y\)，写

\[
  x=\mu_y+v,
  \qquad \|v\|\le r.
\]

比较真实类别和任意 \(j\ne y\)：

\[
\begin{aligned}
  s_y(x)-s_j(x)
  &=
  \langle \mu_y-\mu_j,x\rangle
  -\frac12(\|\mu_y\|^2-\|\mu_j\|^2)\\
  &=
  \frac12\|\mu_y-\mu_j\|^2
  +
  \langle \mu_y-\mu_j,v\rangle.
\end{aligned}
\]

由 Cauchy 不等式和中心分离，

\[
  \langle \mu_y-\mu_j,v\rangle
  \ge
  -\|\mu_y-\mu_j\|\,\|v\|
  \ge
  -r\|\mu_y-\mu_j\|.
\]

令 \(d_{yj}=\|\mu_y-\mu_j\|\ge4r\)，则

\[
  s_y(x)-s_j(x)
  \ge
  \frac12 d_{yj}^2-rd_{yj}
  =
  d_{yj}\left(\frac{d_{yj}}2-r\right)
  \ge
  4r(2r-r)=4r^2\ge 1.
\]

因此对所有 \(j\ne y\)，真实类别 score 至少大 \(1\)，multiclass hinge loss

\[
  \max_{j\ne y}[1+s_j(x)-s_y(x)]_+
\]

为 \(0\)。

\begin{examproblem}[A5：ReLU 网络与 sigmoid 网络表示同一函数]
若有限 ReLU 网络和有限 sigmoid 网络在 \(\R^d\) 上表示同一标量函数，证明该函数必须为常数。
\source{官方卷：2026 Spring, p.~3}
\officialenglish{Let \(d\in\mathbb N\). Consider two feedforward neural networks with the same input dimension \(d\) and scalar output. \(\Phi_{\mathrm{ReLU}}\) uses ReLU in every hidden layer, while \(\Phi_\sigma\) uses sigmoid. Suppose \(\Phi_{\mathrm{ReLU}}(x)=\Phi_\sigma(x)\) for all \(x\in\mathbb R^d\). Show that the represented function must be constant.}
\chinesetranslation{设两个输入维数同为 \(d\)、标量输出的前馈网络，一个所有隐藏层使用 ReLU，另一个使用 sigmoid；若二者在整个 \(\mathbb R^d\) 上表示同一函数，证明该函数只能是常数。}
\end{examproblem}

\solution
ReLU 网络表示的函数是分片仿射（piecewise affine）。Sigmoid 是 real analytic，有限层 sigmoid 网络在常见无直接输入 skip 的前馈结构下表示 real analytic 且有界的函数。

若两者在整个 \(\R^d\) 上相等，则该函数既是 real analytic 又是分片仿射。若分片仿射函数在某个分界超平面两侧有不同 affine piece，则其梯度有跳变，不可能 real analytic。因此所有 piece 的 affine 表达必须拼成同一个全局 affine 函数：

\[
  f(x)=a^\top x+b.
\]

另一方面，sigmoid 网络输出有界，所以全局 affine 函数必须满足 \(a=0\)。因此 \(f(x)=b\) 为常数。

\examnote
若允许 sigmoid 网络含未受限的线性 skip 输出，结论需额外限定网络结构。考试语境通常默认 hidden sigmoid 前馈网络，输出为前层 bounded activations 的线性组合。

\section{B. Deep Learning and Reinforcement Learning}

\begin{examproblem}[B1：backprop、sigmoid/ReLU、momentum、attention scaling、Transformer vs RNN]
回答五个两分小题。
\source{官方卷：2026 Spring, p.~3--4}
\officialenglish{(a) For an MLP \(x_\ell=f_\ell(x_{\ell-1},\theta_\ell)\), \(\ell=1,\ldots,L\), with loss \(J=\mathcal L(x_L,y)\), derive \(\partial J/\partial\theta_\ell\) by the chain rule and state why all gradients by backprop have the order of one forward evaluation. (b) Explain mathematically why \(\sigma(z)=(1+e^{-z})^{-1}\) causes vanishing gradients; contrast ReLU. (c) Write SGD-with-momentum updates for \(v_{k+1}\) and \(x_{k+1}\), and explain their benefit on oscillatory landscapes. (d) For \(W=\operatorname{softmax}(QK^\top/\sqrt{d_k})\), explain \(1/\sqrt{d_k}\) and what fails without it when \(d_k\) is large. (e) Compare RNN and Transformer computational structure and explain why the latter is typically more efficient on modern hardware.}
\chinesetranslation{本题分五问：推导 MLP 的反向传播并解释其复杂度；说明 sigmoid 的梯度消失与 ReLU 的缓解机制；写带动量 SGD 更新并解释其抑制振荡的作用；解释注意力缩放因子 \(1/\sqrt{d_k}\)；比较 RNN 与 Transformer 的计算结构及硬件效率。}
\end{examproblem}

\solution
(a) 对 \(x_\ell=f_\ell(x_{\ell-1},\theta_\ell)\)、\(J=\Lcal(x_L,y)\)，链式法则给出

\[
  \frac{\partial J}{\partial \theta_\ell}
  =
  \frac{\partial J}{\partial x_L}
  \frac{\partial x_L}{\partial x_{L-1}}
  \cdots
  \frac{\partial x_{\ell+1}}{\partial x_\ell}
  \frac{\partial x_\ell}{\partial \theta_\ell}.
\]

\conceptref{concept:jacobian-adjoint}{反向传播}复用中间 adjoint \(\delta_\ell=\partial J/\partial x_\ell\)，每层只做一次局部 Jacobian-vector product，因此总复杂度与前向同阶，差一个常数因子。

(b) Sigmoid 导数

\[
  \sigma'(z)=\sigma(z)(1-\sigma(z))\le 1/4.
\]

深层链式乘积会指数衰减；在饱和区导数更接近 \(0\)。ReLU 导数在 \(z>0\) 为 \(1\)，激活路径上梯度不被系统压小，但 \(z<0\) 时导数为 \(0\)，可能 dead ReLU。

(c) Momentum SGD：

\[
  v_{k+1}=\beta v_k+\nabla f(x_k),
  \qquad
  x_{k+1}=x_k-\eta v_{k+1}.
\]

随机版本用 minibatch gradient。Momentum 累积一致方向、抵消高曲率方向来回震荡，因此在峡谷形 loss landscape 中更快。

(d) 若 \(q,k\) 各维方差约为 \(1\)，则 \(q^\top k\) 方差约为 \(d_k\)。除以 \(\sqrt{d_k}\) 使 logits 方差稳定。若去掉缩放，大 \(d_k\) 时 softmax 饱和到近 one-hot，梯度变小，训练不稳定。

(e) RNN 必须按时间递推，序列长度 \(L\) 的依赖路径长且难并行。Transformer 每层 \conceptref{concept:self-attention}{self-attention} 可对所有 token 并行计算，任意 token 间一层可交互，更适合 GPU/TPU 批量矩阵乘。

\begin{examproblem}[B2：Flow matching 与 stochastic interpolant]
从 base distribution \(\pi_0\) 到 data distribution \(\pi_1\)，用线性插值 \(x_t=(1-t)x_0+tx_1\) 定义速度场，证明 continuity equation、回归目标和带 score 的 SDE 有同样 marginal。
\source{官方卷：2026 Spring, p.~4--5}
\officialenglish{Let \(\pi_0\) be a simple base distribution on \(\mathbb R^d\) and \(\pi_1\) a data distribution. Seek \(v_t\) such that \(\frac d{dt}x_t=v_t(x_t)\), \(x_0\sim\pi_0\), transports \(\pi_0\) to \(\pi_1\) at \(t=1\). (a) For coupling \((x_0,x_1)\sim\gamma\) and \(x_t=(1-t)x_0+tx_1\), compute \(v_t(x_t\mid x_0,x_1)\). (b) Define \(v_t(x)=\mathbb E[v_t(x_t\mid x_0,x_1)\mid x_t=x]\); show its flow satisfies \(\partial_t\pi_t+\nabla\cdot(\pi_tv_t)=0\). (c) Give a squared regression objective for \(v_\theta(t,x)\) and prove its minimizer is this marginal velocity. (d) Derive the Fokker--Planck equation of \(dx_t=(v_t(x_t)+\beta_t\nabla_x\log\pi_t(x_t))dt+\sqrt{2\beta_t}\,dW_t\), show it has the same marginals, and explain score matching for the unknown score.}
\chinesetranslation{从基分布 \(\pi_0\) 到数据分布 \(\pi_1\) 构造流匹配向量场。先对线性插值得到条件速度，再定义条件期望速度并证明连续性方程；用平方回归学习参数化速度；最后给出加入扩散项的 SDE 的 Fokker--Planck 方程、同边缘分布结论，并说明如何用 score matching 估计未知 score。}
\end{examproblem}

\solution
(a) 条件速度为

\[
  v_t(x_t|x_0,x_1)
  =
  \frac{d}{dt}x_t
  =
  x_1-x_0.
\]

(b) 定义 marginal velocity

\[
  v_t(x)=
  \E[x_1-x_0\mid x_t=x].
\]

对任意 test function \(\varphi\)，

\[
  \frac{d}{dt}\E\varphi(x_t)
  =
  \E\langle \nabla\varphi(x_t),x_1-x_0\rangle
  =
  \E\langle \nabla\varphi(x_t),v_t(x_t)\rangle.
\]

若 \(\pi_t\) 是 \(x_t\) 的密度，则

\[
  \frac{d}{dt}\int \varphi(x)\pi_t(x)\,dx
  =
  \int \langle \nabla\varphi(x),v_t(x)\rangle\pi_t(x)\,dx.
\]

分部积分得到弱形式

\[
  \partial_t\pi_t+\nabla\cdot(\pi_t v_t)=0.
\]

因此 ODE \(dx_t/dt=v_t(x_t)\) 产生同一组 marginals，从 \(\pi_0\) transport 到 \(\pi_1\)。

(c) 用模型 \(v_\theta(t,x)\) 最小化

\[
  \mathcal{L}(\theta)
  =
  \E_{t,(x_0,x_1),x_t}
  \|v_\theta(t,x_t)-(x_1-x_0)\|^2.
\]

对固定 \((t,x)\)，平方损失的最优预测是条件均值：

\[
  v_\theta^*(t,x)
  =
  \E[x_1-x_0\mid x_t=x]
  =
  v_t(x).
\]

(d) 考虑 SDE

\[
  dx_t=
  \left(v_t(x_t)+\beta_t\nabla_x\log\pi_t(x_t)\right)dt
  +
  \sqrt{2\beta_t}\,dW_t.
\]

Fokker-Planck equation 为

\[
  \partial_t p_t
  =
  -\nabla\cdot\left[
    \left(v_t+\beta_t\nabla\log\pi_t\right)p_t
  \right]
  +
  \beta_t\Delta p_t.
\]

若 \(p_t=\pi_t\)，则

\[
  -\nabla\cdot(\beta_t\pi_t\nabla\log\pi_t)+\beta_t\Delta\pi_t
  =
  -\beta_t\nabla\cdot(\nabla\pi_t)+\beta_t\Delta\pi_t
  =
  0.
\]

于是

\[
  \partial_t\pi_t=-\nabla\cdot(\pi_t v_t),
\]

与 deterministic ODE 的 continuity equation 一致，故 marginals 相同。实际中 \(\nabla\log\pi_t\) 未知，可用 \conceptref{concept:diffusion-score}{denoising score matching} 训练 \(s_\eta(t,x)\approx \nabla\log\pi_t(x)\)。

\begin{examproblem}[B3：Bellman、policy iteration、actor-critic、RLHF/DPO]
回答强化学习基础与偏好对齐算法。
\source{官方卷：2026 Spring, p.~5--6}
\officialenglish{For discounted MDP \((\mathcal S,\mathcal A,P,r,\gamma)\): (a) derive the Bellman equation for \(V^\pi(s)=\mathbb E_\pi[\sum_{t\ge0}\gamma^tr(s_t,a_t)\mid s_0=s]\); (b) state policy iteration, including evaluation and improvement; (c) when \(P\) is unknown, describe actor--critic replacements; (d) for human preferences \((x,y^+,y^-)\), specify one RLHF update objective for \(\pi_\theta(y\mid x)\), either reward-model RLHF with reference-policy KL or DPO-style direct preference optimisation, defining your notation.}
\chinesetranslation{对折扣 MDP：推导价值函数的 Bellman 方程；写清策略迭代的评估与改进；在模型未知时说明 actor--critic 如何替代两步；再用偏好三元组给出一种 RLHF 的策略更新目标，可选择带 KL 正则的奖励模型 RLHF 或 DPO。}
\end{examproblem}

\solution
(a) \conceptref{concept:bellman}{Bellman equation}：

\[
  V^\pi(s)
  =
  \sum_a\pi(a|s)
  \left[
    r(s,a)+\gamma\sum_{s'}P(s'|s,a)V^\pi(s')
  \right].
\]

(b) Policy iteration：

\[
  V^{\pi_k}
  =
  T^{\pi_k}V^{\pi_k}
  \quad\text{(policy evaluation)}
\]

然后

\[
  \pi_{k+1}(s)\in
  \argmax_a
  \left[
    r(s,a)+\gamma\sum_{s'}P(s'|s,a)V^{\pi_k}(s')
  \right]
  \quad\text{(policy improvement)}.
\]

(c) Model-free actor-critic 不显式知道 \(P\)。critic 用采样 TD error

\[
  \delta_t=r_t+\gamma V_\phi(s_{t+1})-V_\phi(s_t)
\]

学习 \(V_\phi\) 或 \(Q_\phi\)；actor 用

\[
  \nabla_\theta\log\pi_\theta(a_t|s_t)\,\widehat A_t
\]

更新策略，其中 \(\widehat A_t\) 可由 TD error 或 GAE 估计；policy gradient 的 log-derivative 推导回查 \conceptref{concept:policy-gradient}{第三章策略梯度}。

(d) Reward-model-based \conceptref{concept:rlhf-dpo}{RLHF}：先从偏好数据 \((x,y^+,y^-)\) 训练

\[
  \mathcal{L}_{\mathrm{RM}}(\psi)
  =
  -\log\sigma(r_\psi(x,y^+)-r_\psi(x,y^-)).
\]

再优化

\[
  \max_\theta
  \E_{x,y\sim\pi_\theta}
  \left[
    r_\psi(x,y)
    -
    \beta
    \log\frac{\pi_\theta(y|x)}{\pi_{\mathrm{ref}}(y|x)}
  \right].
\]

\conceptref{concept:rlhf-dpo}{DPO} 直接优化

\[
  \mathcal{L}_{\mathrm{DPO}}(\theta)
  =
  -
  \E
  \log
  \sigma
  \left[
    \beta
    \left(
      \log\frac{\pi_\theta(y^+|x)}{\pi_{\mathrm{ref}}(y^+|x)}
      -
      \log\frac{\pi_\theta(y^-|x)}{\pi_{\mathrm{ref}}(y^-|x)}
    \right)
  \right].
\]

\section{C. Optimization Methods in Artificial Intelligence}

\begin{examproblem}[C1-C10：prox-SGD、非扩张、噪声地板与平均]
复合优化

\[
  \min_x f(x)+R(x)
\]

下，证明 proximal optimality、fixed point、monotonicity、nonexpansiveness、one-step inequality、variance decomposition、linear contraction、averaging variance 和 averaged convergence。
\source{官方卷：2026 Spring, p.~6--8}
\officialenglish{Consider \(\min_{x\in\mathbb R^d}F(x)=f(x)+R(x)\), where \(f\) is differentiable, \(L\)-smooth and \(\mu\)-strongly convex, and \(R\) is proper, closed and convex. Let \(g(x,\xi)=\nabla f(x)+\xi\), \(\mathbb E\xi=0\), \(\mathbb E\|\xi\|^2\le\sigma^2\), and \(x_{k+1}=\operatorname{prox}_{\gamma R}(x_k-\gamma g_k)\); set \(\bar x^k=k^{-1}\sum_{t=1}^kx^t\). (1) Define smoothness and strong convexity. (2) Prove \(x^+=\operatorname{prox}_{\gamma R}(u)\iff\gamma^{-1}(u-x^+)\in\partial R(x^+)\). (3) Show the unique minimizer \(x^*\) is a proximal fixed point. (4) Prove monotonicity of \(\partial R\). (5) Prove nonexpansiveness of prox. (6) Establish the one-step inequality. (7) Take conditional expectation and decompose variance. (8) For \(\gamma=2/(L+\mu)\), prove linear contraction with noise floor and interpret it. (9) Under uncorrelated errors of bounded second moment \(V\), show \(\mathbb E\|\bar e^k\|^2\le V/k\). (10) With \(R\equiv0\), derive the smoothness one-step bound and, for \(\gamma_t=1/(\mu t)\), show \(\mathbb E[f(\bar x^k)-f^*]=O(\log k/k)\).}
\chinesetranslation{复合强凸优化的完整十问：定义；prox 的最优性条件、固定点、次微分单调性和非扩张；prox-SGD 的一步界、条件期望与噪声地板；平均迭代的方差缩减；以及 \(R=0\) 时用 \(\gamma_t=1/(\mu t)\) 得到 \(O(\log k/k)\) 的平均收敛率。所有精确常数与提示见官方第 6--8 页。}
\end{examproblem}

\solution
\textbf{C1-C5。} \conceptref{concept:prox}{prox 定义}和 \conceptref{concept:prox-nonexpansive}{prox 非扩张推导}见第 \ref{ch:primer} 章。prox 最优性来自 convex problem 一阶条件：

\[
  0\in \partial R(x^+)+\frac1\gamma(x^+-u)
  \quad\Longleftrightarrow\quad
  \frac{u-x^+}{\gamma}\in\partial R(x^+).
\]

最优点 \(x^*\) 满足

\[
  0\in \nabla f(x^*)+\partial R(x^*),
\]

所以

\[
  \frac{x^*-\gamma\nabla f(x^*)-x^*}{\gamma}
  =
  -\nabla f(x^*)\in\partial R(x^*),
\]

即

\[
  x^*=\prox_{\gamma R}(x^*-\gamma\nabla f(x^*)).
\]

若 \(r_x\in\partial R(x)\)，\(r_y\in\partial R(y)\)，convex subgradient inequality 给出

\[
  R(y)\ge R(x)+\langle r_x,y-x\rangle,
  \qquad
  R(x)\ge R(y)+\langle r_y,x-y\rangle.
\]

相加得

\[
  \langle r_x-r_y,x-y\rangle\ge0.
\]

设 \(x^+=\prox_{\gamma R}(u)\)，\(y^+=\prox_{\gamma R}(v)\)。由最优性，

\[
  \frac{u-x^+}{\gamma}\in\partial R(x^+),
  \qquad
  \frac{v-y^+}{\gamma}\in\partial R(y^+).
\]

单调性给出

\[
  \langle u-v-(x^+-y^+),x^+-y^+\rangle\ge0.
\]

因此

\[
  \|x^+-y^+\|^2\le \langle u-v,x^+-y^+\rangle
  \le \|u-v\|\,\|x^+-y^+\|,
\]

即 \conceptref{concept:prox-nonexpansive}{prox 非扩张}。

\textbf{C6-C8。} 用 fixed point 和非扩张：

\[
  \|x_{k+1}-x^*\|^2
  \le
  \|x_k-x^*-\gamma(g^k-\nabla f(x^*))\|^2.
\]

条件期望下，\(g^k=\nabla f(x_k)+\xi^k\)，\(\E_k\xi^k=0\)，得

\[
  \E_k\|x_{k+1}-x^*\|^2
  \le
  \|x_k-x^*-\gamma(\nabla f(x_k)-\nabla f(x^*))\|^2+\gamma^2\sigma^2.
\]

取 \(\gamma=2/(L+\mu)\) 并用强凸-光滑联合不等式，得到

\[
  \E\|x_{k+1}-x^*\|^2
  \le
  (1-\rho)\E\|x_k-x^*\|^2+\gamma^2\sigma^2,
  \qquad
  \rho=\frac{4\mu L}{(L+\mu)^2}.
\]

递推极限噪声地板为

\[
  \frac{\gamma^2\sigma^2}{\rho}.
\]

含义是：即使均值方向收缩，随机梯度每步注入方差，常数步长下误差不会收敛到零。

\textbf{C9。} 令 \(\bar e_k=k^{-1}\sum_{t=1}^{k}e_t\)。若 \(e_t\) 两两不相关且 \(\E\|e_t\|^2\le V\)，则

\[
\begin{aligned}
  \E\|\bar e_k\|^2
  &=
  \E\left\|
    \frac1k\sum_{t=1}^{k}e_t
  \right\|^2\\
  &=
  \frac1{k^2}\sum_{t=1}^{k}\E\|e_t\|^2
  +
  \frac{2}{k^2}\sum_{s<t}\E\langle e_s,e_t\rangle\\
  &\le \frac{V}{k}.
\end{aligned}
\]

这解释了平均迭代比最后迭代方差小。

\textbf{C10。} 当 \(R\equiv0\) 时，由 C7：

\[
  \E_t\|x_{t+1}-x^*\|^2
  \le
  \|x_t-x^*\|^2
  -2\gamma_t\langle\nabla f(x_t),x_t-x^*\rangle
  +\gamma_t^2\|\nabla f(x_t)\|^2
  +\gamma_t^2\sigma^2.
\]

凸性给出

\[
  \langle\nabla f(x_t),x_t-x^*\rangle\ge f(x_t)-f^*.
\]

光滑性给出

\[
  \|\nabla f(x_t)\|^2\le 2L(f(x_t)-f^*).
\]

因此

\[
  \E_t\|x_{t+1}-x^*\|^2
  \le
  \|x_t-x^*\|^2
  -
  2\gamma_t(1-L\gamma_t)(f(x_t)-f^*)
  +
  \gamma_t^2\sigma^2.
\]

若 \(\gamma_t\le 1/(2L)\)，则

\[
  \E_t\|x_{t+1}-x^*\|^2
  \le
  \|x_t-x^*\|^2
  -
  \gamma_t(f(x_t)-f^*)
  +
  \gamma_t^2\sigma^2.
\]

对 \(t=1,\ldots,k\) 求和并取 \(\gamma_t=1/(\mu t)\)，左侧 telescoping，右侧噪声和为

\[
  \sum_{t=1}^{k}\gamma_t^2\sigma^2=O(1).
\]

利用强凸情形下加权平均或题目简化提示 \(\sum_{t=1}^{k}1/t\sim\log k\)，可得

\[
  \E[f(\bar x^k)-f^*]
  \le
  O\left(\frac{\log k}{k}\right).
\]

\section{D. Natural Language Processing}

\begin{examproblem}[D1：InfoNCE 梯度与高维负样本]
给定单位向量 \(z,z^+,z_i^-\) 和

\[
  \Lcal=-\log\frac{\exp(\langle z,z^+\rangle/\tau)}
  {\exp(\langle z,z^+\rangle/\tau)+\sum_i\exp(\langle z,z_i^-\rangle/\tau)},
\]

求 \(\nabla_z\Lcal\)，并解释高维负样本相似度尺度和 false negatives。
\source{官方卷：2026 Spring, p.~8}
\officialenglish{An encoder \(f_\theta\) maps sentence \(x\) to a unit vector \(z\in\mathbb R^d\). Given positive \(z^+\), negatives \(z_i^-\), and temperature \(\tau\), consider \(\mathcal L=-\log\frac{\exp(\langle z,z^+\rangle/\tau)}{\exp(\langle z,z^+\rangle/\tau)+\sum_{i=1}^n\exp(\langle z,z_i^-\rangle/\tau)}\). (a) Compute \(\nabla_z\mathcal L\) and interpret pull/push on the unit hypersphere. (b) Under isotropic negatives in high dimension, give the scale of \(\langle z,z_i^-\rangle\) and its maximum over \(n\) negatives, then explain the damage caused by false negatives.}
\chinesetranslation{对单位球上的 InfoNCE：求锚向量 \(z\) 的梯度并作“拉近正样本、推离负样本”的切空间解释；再在高维各向同性假设下估计负样本内积及最大内积的尺度，并解释语义相同的假负样本为何有害。}
\end{examproblem}

\solution
令候选集合 \(u_0=z^+\)，\(u_i=z_i^-\)，logits 为 \(a_j=\langle z,u_j\rangle/\tau\)，softmax 权重为

\[
  p_j=\frac{e^{a_j}}{\sum_{\ell=0}^{n}e^{a_\ell}}.
\]

则

\[
  \Lcal=-a_0+\log\sum_{\ell=0}^{n}e^{a_\ell},
\]

所以

\[
  \nabla_z\Lcal
  =
  \frac1\tau
  \left(
    \sum_{j=0}^{n}p_j u_j-z^+
  \right)
  =
  \frac1\tau
  \left[
    (p_0-1)z^+
    +
    \sum_{i=1}^{n}p_i z_i^-
  \right].
\]

负梯度方向把 \(z\) 拉向 \(z^+\)，推离权重大、相似度高的 negatives。在单位球约束下，实际更新可投影到切空间。

高维单位球上，若 \(z_i^-\) isotropic，则

\[
  \E\langle z,z_i^-\rangle=0,
  \qquad
  \Var(\langle z,z_i^-\rangle)\approx \frac1d.
\]

最大负样本相似度的典型尺度约为

\[
  \max_{i\le n}\langle z,z_i^-\rangle
  \approx
  \sqrt{\frac{2\log n}{d}}.
\]

False negative 若语义上与 anchor 同类，会有异常高相似度，被 loss 强行推开，破坏类内结构并增加训练噪声。

\begin{examproblem}[D2：监督主题模型]
为带优先级标签的客服工单设计 supervised topic model，解释 supervision 改变 topics 的原因，写 label model、sampling process 和 joint。
\source{官方卷：2026 Spring, p.~8--9}
\officialenglish{Customer-support tickets have token words \(w_{dn}\in\{1,\ldots,V\}\) and priority labels \(y_d\in\{0,1\}\). For a supervised topic model with \(K\) topics, \(\theta_d\sim\operatorname{Dir}(\alpha)\), \(\phi_k\sim\operatorname{Dir}(\eta)\), assignments \(z_{dn}\), and empirical topic frequency \(\bar z_d\): (a) explain why observing \(y_d\) changes topics and name a vanilla-LDA failure it mitigates; (b) propose \(p(y_d\mid\bar z_d,\beta)\), e.g. logistic regression, and interpret \(\beta\); (c) specify the conditional sampling process and the factorized joint \(p(w,z,\Theta,\Phi,y\mid\alpha,\eta,\beta)\).}
\chinesetranslation{电商工单包含词 token 与二元优先级标签。对带 Dirichlet 先验的监督主题模型：(a) 说明标签为何会改变主题以及其可修复的普通 LDA 失败模式；(b) 用逻辑回归等给出 \(p(y_d\mid\bar z_d,\beta)\) 并解释 \(\beta\)；(c) 给出条件生成步骤和包含所有文档、主题、token 与标签项的联合分布分解。}
\end{examproblem}

\solution
加入标签 \(y_d\) 会让 topic 不只解释词共现，还要解释预测任务。Vanilla LDA 可能发现“模板词/平台词”等高频主题，却忽略区分 urgent/normal 的低频故障词；supervision 会推动 topic 对标签有判别力。

取 logistic label model：

\[
  \sigma(u)=\frac{1}{1+e^{-u}},
  \qquad
  p(y_d=1|\bar z_d,\beta)=\sigma(\beta^\top \bar z_d),
\]

\[
  p(y_d|\bar z_d,\beta)
  =
  \sigma(\beta^\top\bar z_d)^{y_d}
  \left(1-\sigma(\beta^\top\bar z_d)\right)^{1-y_d}.
\]

\(\beta_k\) 表示第 \(k\) 个 topic 对 urgent log-odds 的贡献。

采样过程：

\[
  \theta_d\sim\Dir(\alpha),
  \qquad
  \phi_k\sim\Dir(\eta).
\]

对每个 token：

\[
  z_{dn}\sim\Cat(\theta_d),
  \qquad
  w_{dn}\sim\Cat(\phi_{z_{dn}}).
\]

计算 \(\bar z_d\)，再采样

\[
  y_d\sim\Bern(\sigma(\beta^\top \bar z_d)).
\]

joint 为

\[
  p(w,z,\Theta,\Phi,y|\alpha,\eta,\beta)
  =
  \prod_k p(\phi_k|\eta)
  \prod_d
  p(\theta_d|\alpha)
  \left[
    \prod_{n=1}^{N_d}
    p(z_{dn}|\theta_d)p(w_{dn}|\phi_{z_{dn}})
  \right]
  p(y_d|\bar z_d,\beta).
\]

\begin{examproblem}[D3：长上下文 Transformer 设计]
设计可处理 \(L\gg10^4\) 的 Transformer variant，保留局部和长程依赖，并分析复杂度。
\source{官方卷：2026 Spring, p.~9}
\officialenglish{Standard self-attention costs \(O(L^2)\). Design a Transformer variant handling \(L\gg10^4\) while preserving local and long-range dependencies. Choose block-sparse attention, sliding-window plus global tokens, low-rank/linear attention, or recurrent memory. (a) Precisely define its computation/pattern and hyperparameters. (b) Give per-layer dominant time and memory complexity in \(L,d\) and those hyperparameters. (c) Give an argument that long-range interactions remain expressible.}
\chinesetranslation{标准自注意力为 \(O(L^2)\)。任选块稀疏、滑动窗加全局 token、低秩/线性注意力或循环记忆，设计可处理 \(L\gg10^4\) 的 Transformer：(a) 精确定义注意力模式与超参数；(b) 分析每层时间与内存；(c) 论证仍可表达长程相互作用。}
\end{examproblem}

\solution
本题是在 \conceptref{concept:self-attention}{self-attention} 基础上修改可见集合，因此先回查第三章的 \(Q,K,V\)、缩放因子和复杂度定义。

选择 sliding window + global tokens。设窗口半径为 \(w\)，global token 集合 \(G\)，\(|G|=g\)。对普通位置 \(i\)，定义可见集合

\[
  \mathcal{N}(i)=\{j:|i-j|\le w\}\cup G.
\]

对 global token \(g_a\)，令

\[
  \mathcal{N}(g_a)=\{1,\ldots,L\}\cup G.
\]

Attention 只在 \(\mathcal{N}(i)\) 内计算：

\[
  \operatorname{Attn}(i)
  =
  \sum_{j\in \mathcal{N}(i)}
  \softmax_{j}
  \left(
    \frac{q_i^\top k_j}{\sqrt d}
  \right)v_j.
\]

普通 token 的成本是 \(O(L(w+g)d)\)，global token 成本是 \(O(gLd)\)。总 attention 时间和显存主项为

\[
  O(L(w+g)d+gLd)=O(L(w+g)d).
\]

若 \(w,g\ll L\)，比标准 \(O(L^2d)\) 低。长程信息通过 global tokens 两跳传播：位置 \(i\to G\to j\)，一层可把所有 token 写入 global，下一层可从 global 读出。因此 attention graph 直径为常数级，仍可表达长程依赖。

\begin{examproblem}[D4：信息抽取的 MLN 全局一致性]
给定实体链接和关系候选的局部置信号，用 MLN 约束全局一致性，并处理同名人物跨时代歧义。
\source{官方卷：2026 Spring, p.~9--10}
\officialenglish{An information-extraction pipeline gives uncertain mention--entity and relation candidates. An MLN uses unknown \(\operatorname{Link}(m,e)\), \(\operatorname{Rel}(e_1,e_2,r)\) and evidence \(\operatorname{Cand},\operatorname{Type},\operatorname{HighLink},\operatorname{HighRel}\). (a) Write weighted formulas propagating high-confidence local evidence, enforcing at-most-one entity per mention, and relation-type consistency; discuss hard/soft weights. (b) Describe concrete upstream models and objectives producing \(\operatorname{HighLink}\) and \(\operatorname{HighRel}\). (c) Extend the MLN for two same-name persons 500 years apart, adding predicates/rules for relation-induced contemporaneity and birth-year incompatibility, and give a minimal joint-inference example.}
\chinesetranslation{构建新闻信息抽取的 MLN：未知谓词是实体链接与关系，观测证据包括候选、类型及高置信局部信号。(a) 写带权规则，处理证据传播、每个 mention 至多一个实体、关系类型一致性，并区分硬软约束；(b) 说明上游如何建模并产生高置信信号；(c) 对同名但相差五百年的人加入时代一致性和出生年冲突规则，给一个最小联合推断例子。}
\end{examproblem}

\solution
可用以下 soft/hard rules：

\[
  \operatorname{HighLink}(m,e)\Rightarrow \operatorname{Link}(m,e)
  \quad\text{(soft positive, weight high)}.
\]

\[
  \operatorname{HighRel}(m_1,m_2,r)\wedge
  \operatorname{Link}(m_1,e_1)\wedge
  \operatorname{Link}(m_2,e_2)
  \Rightarrow
  \operatorname{Rel}(e_1,e_2,r)
  \quad\text{(soft)}.
\]

每个 mention 至多链接一个 entity：

\[
  \operatorname{Link}(m,e_1)\wedge \operatorname{Link}(m,e_2)\Rightarrow e_1=e_2
  \quad\text{(hard)}.
\]

类型一致性：

\[
  \operatorname{Rel}(e_1,e_2,r)\wedge \operatorname{Requires}(r,t_1,t_2)
  \Rightarrow
  \operatorname{Type}(e_1,t_1)\wedge \operatorname{Type}(e_2,t_2).
\]

HighLink 可由 bi-encoder 或 cross-encoder entity linking model 产生：输入 mention context 与 candidate entity description，输出 \(p(e|m)\)，通过阈值、top-\(k\) 或 margin 得到 evidence。HighRel 可由 relation classifier 产生：输入两个 mentions 的上下文、entity types 和 dependency/\conceptref{concept:self-attention}{Transformer span representation}，输出 \(p(r|m_1,m_2)\)。

同名跨时代扩展谓词：\(\operatorname{BirthYear}(e,b)\)、\(\operatorname{SameEra}(e_1,e_2)\)、\(\operatorname{EraRel}(r)\)。规则：

\[
  \operatorname{HighRel}(m_1,m_2,r)\wedge
  \operatorname{Link}(m_1,e_1)\wedge
  \operatorname{Link}(m_2,e_2)\wedge
  \operatorname{EraRel}(r)
  \Rightarrow
  \operatorname{SameEra}(e_1,e_2).
\]

\[
  \operatorname{BirthYear}(e_1,b_1)\wedge
  \operatorname{BirthYear}(e_2,b_2)\wedge
  |b_1-b_2|>120
  \Rightarrow
  \neg\operatorname{SameEra}(e_1,e_2).
\]

第二条可设 hard 或极高负权。例：文档出现“张三与李四合作发表论文”，候选张三有古代人物 \(e_{\mathrm{old}}\) 和现代学者 \(e_{\mathrm{new}}\)，李四为现代学者，关系 \(\operatorname{Coauthor}\) 是 EraRel。HighRel 支持合作关系会诱导 SameEra；BirthYear 差 500 年惩罚 \(e_{\mathrm{old}}\)，从而联合推断偏向 \(e_{\mathrm{new}}\)。

\begin{examproblem}[D5：RAG memory gain 与 MDP 优化]
定义 memory system 的 per-turn gain，推导 \(\Delta>0\) 条件，并把 memory 写入、检索、淘汰机制表述为 MDP。
\source{官方卷：2026 Spring, p.~10--11}
\officialenglish{A dialog NLU/QA system uses RAG memory: it writes \(c_t=\operatorname{Memorize}(H_t)\) to \(M_t\), retrieves \(R_t=\operatorname{Retrieve}(H_t,M_t,k)\), and answers \(\hat y_t=f_\theta(H_t,R_t)\). (a) With \(\Delta_t=s(f_\theta(H_t,R_t),y_t^*)-s(f_\theta(H_t,\varnothing),y_t^*)\), retrieval failure probability \(\varepsilon\), and conditional gains \(g_c,g_{\neg c}\), derive the necessary and sufficient condition for total gain \(\Delta>0\); when \(g_c>g_{\neg c}\), show \(\varepsilon<g_c/(g_c-g_{\neg c})\). (b) Formulate memory management as an MDP and RL objective. State state, action including write/merge/evict and retrieval choices, and a quality--cost reward.}
\chinesetranslation{对话式 RAG 外部记忆系统先写入、再检索、后生成。(a) 从检索失败概率 \(\varepsilon\) 与“检到关键记忆/未检到”两种条件收益导出总增益为正的充要条件，并给出题中阈值；(b) 将写入、合并、淘汰以及检索参数共同建模为 MDP 动作，给出状态、奖励和 RL 优化目标。}
\end{examproblem}

\solution
本题的 memory/RAG 系统分解回查 \conceptref{concept:rag-mln}{第三章 embedding 知识库推理与 RAG}，把写入、检索、淘汰动作形式化为 MDP 时回查 \conceptref{concept:mdp}{MDP 定义}。

由全期望公式，

\[
  \Delta
  =
  \E[\Delta_t]
  =
  (1-\varepsilon)g_c+\varepsilon g_{\neg c}.
\]

因此

\[
  \Delta>0
  \quad\Longleftrightarrow\quad
  (1-\varepsilon)g_c+\varepsilon g_{\neg c}>0.
\]

若 \(g_c>g_{\neg c}\)，整理：

\[
  g_c-\varepsilon(g_c-g_{\neg c})>0
  \quad\Longleftrightarrow\quad
  \varepsilon<\frac{g_c}{g_c-g_{\neg c}}.
\]

当 \(\varepsilon\) 大，即经常检索不到关键记忆时，系统收益由 \(g_{\neg c}\) 主导；若错误记忆还会误导生成，\(g_{\neg c}\) 可能为负，稳定正收益很难保证。

MDP 表述：

\[
  s_t=(H_t,M_t,\text{budget state},\text{retrieval diagnostics},\text{last answer feedback}).
\]

动作 \(a_t\) 包括两类：

\begin{itemize}
  \item 写入侧：write/merge/evict、摘要压缩粒度、tag schema、TTL、去重阈值；
  \item 检索侧：\(k\)、dense/sparse hybrid weights、query rewrite、reranker depth、是否引用长期记忆。
\end{itemize}

奖励可设

\[
  r_t=
  s(\hat y_t,y_t^*)-
  s(\hat y_t^{\emptyset},y_t^*)
  -
  \lambda\,\operatorname{Cost}(a_t)
  -
  \eta\,\operatorname{Staleness}(R_t).
\]

目标是

\[
  \max_\pi
  \E_\pi
  \sum_{t=0}^{T}\gamma^t r_t.
\]

实际训练可用离线日志做 contextual bandit/IPS 评估，或在模拟环境中用 actor-critic 优化记忆策略；上线时需要 budget hard constraint 和回滚策略，避免写入污染长期记忆。
